\StartChapter{Systeme von Differentialgleichungen}

\SubsectionBar{Eigenschaften}

\textVorBox{Treten zwei Differentialgleichungen gekoppelt auf, so spricht man von einem \ZSFkeyword{System von Differentialgleichungen 1. Ordnung}:}
\begin{formulabox}
    $\displaystyle
    \begin{cases}
        y_1' = f_1(x,y_1,y_2,b_1)\\
        y_2' = f_2(x,y_1,y_2,b_2)
    \end{cases}$
\end{formulabox}
\textNachBox{Dabei sind $y_1$ und $y_2$ Funktionen in $x$.}

Wenn $f_1$ und $f_2$ in $x$ stetig und nach $y_1$ und $y_2$ stetig differenzierbar sind, dann hat das AWP $y_1(x_0)=y_{1,0}$, $y_2(x_0)=y_{2,0}$ eine Lösung.

\textVorBox{Ein \ZSFkeyword{autonomes System} ist von der Form:}
\begin{formulabox}
    $\displaystyle
    \begin{cases}
        y_1' = f_1(y_1,y_2,b_1)\\
        y_2' = f_2(y_1,y_2,b_2)
    \end{cases}$
\end{formulabox}
\textNachBox{Dabei erfolgt die einzige Abhängigkeit von der Variablen $x$ indirekt über $y_1(x)$ und $y_2(x)$.}

Eine \ZSFkeyword{Trajektorie} ist die durch $(x_0,y_0)$ gehende Feldlinie. Die Schar aller Trajektorien ist das \ZSFkeyword{Phasenportrait}, ein Vektorfeld:
\begin{formulabox}
    $\displaystyle
    \begin{cases}
        x' = f_1(x,y)\\
        y' = f_2(x,y)
    \end{cases}
    \Rightarrow
    \vec v=
    \begin{pmatrix}
        x'\\
        y'
    \end{pmatrix}
    =
    \begin{pmatrix}
        v_1\\
        v_2
    \end{pmatrix}$
\end{formulabox}

\SubsectionBar{Linear-autonome DGL mit konst. Koeffizienten}

\textVorBox{Ein System der folgenden Form ist ein \ZSFkeyword{linear autonomes DGL-System mit konstanten Koeffizienten} der Ordnung 1:}
\begin{formulabox}
    $\displaystyle
    \begin{cases}
        x_1' = a_{11}x_1 + a_{12}x_2 + b_1\\
        x_2' = a_{21}x_1 + a_{22}x_2 + b_2
    \end{cases}$
\end{formulabox}

\begin{procedure}[Lösungsverfahren: Homogene DGL-Systeme]
    \ProcStep{System aufstellen}{Matrix $A$ und Vektoren $\dot{\vec x}$, $\vec x$, $\vec b$ aufstellen: $\dot{\vec x}=A\vec x$.}
    \ProcStep{Eigenwerte/-vektoren}{Eigenwerte $\lambda_i$ und Eigenvektoren $\vec v_i$ bestimmen.}
    \ProcStep{Homogene Lösung}{$\vec x=\sum_i C_i\cdot e^{\lambda_i t}\cdot \vec v_i$.}
\end{procedure}

\textVorBox{Bei einem imaginären Eigenwertepaar $\lambda_{1,2}=a\pm bi$ gilt in reeller Darstellung:}
\begin{formulabox}
    $\displaystyle
    \vec x=
    C_1\cdot e^{at}
    \begin{pmatrix}
        c\cos(bt)\mp d\sin(bt)\\
        e\cos(bt)\mp f\sin(bt)\\
        g\cos(bt)\mp h\sin(bt)
    \end{pmatrix}
    +C_2\cdot e^{at}
    \begin{pmatrix}
        c\sin(bt)\pm d\cos(bt)\\
        e\sin(bt)\pm f\cos(bt)\\
        g\sin(bt)\pm h\cos(bt)
    \end{pmatrix}$
\end{formulabox}

\textVorBox{Wenn die Matrix $A$ nicht diagonalisierbar ist oder das DGLS inhomogen ist, eignet sich das \ZSFkeyword{Eliminationsverfahren}:}
\begin{procedure}
    \ProcStep{Zweite Ableitung}{$x_1'$ oder $x_2'$ zu $x_1''$ bzw. $x_2''$ ableiten.}
    \ProcStep{Ersetzen}{$x_1'$ bzw. $x_2'$ im abgeleiteten Term durch die ursprüngliche Gleichung ersetzen.}
    \ProcStep{Eliminieren}{$x_1$ bzw. $x_2$ mit Hilfe der ursprünglichen Gleichung eliminieren.}
    \ProcStep{DGL lösen}{Zu bekannter DGL umformen und Nullstellen bestimmen.}
    \ProcStep{Zweite Funktion}{Aus $x_1$ bzw. $x_2$ die zweite Funktion bestimmen.}
\end{procedure}

\SubsectionBar{Gleichgewichtspunkte}

\textVorBox{Wenn die beiden konstanten Funktionen $x_1(t)=c_1$ und $x_2(t)=c_2$ das autonome System lösen:}
\begin{formulabox}
    $\displaystyle
    \begin{cases}
        x_1' = f_1(x_1,x_2)\\
        x_2' = f_2(x_1,x_2)
    \end{cases}$
\end{formulabox}
\textNachBox{Dann ist $(c_1,c_2)$ ein \ZSFkeyword{Gleichgewichtspunkt} des DGL-Systems.}

\begin{propertylist}[Gleichgewichtspunkt finden]
    \item \textbf{DGL-System:} $x_1'=x_2'=0$ setzen und nach $x_1, x_2$ auflösen.
    \item \textbf{Gewöhnliche DGL:} $y'=f(y(t))$ mit $f(y_g)=0$ --- das AWP $y(t_0)=y_g$ hat die konstante Lösung $y(t)=y_g$.
\end{propertylist}

Die Gleichungssysteme lassen sich umschreiben zu $\dot{\vec x}=A\vec x$. Die Eigenwerte $\lambda_i$ von $A$ geben Aufschluss über die Art des Gleichgewichtspunkts:
\begin{fulltablebox}{Z{0.85} | Z{1.15}}
    \ZSFrowColor
    Asymptotisch stabil: & Alle Eigenwerte von $A$ haben negativen Realteil \\ \hline
    Stabil: & Die Eigenwerte von $A$ haben Realteil $\mathcal{R}(\lambda_i)\leq 0$ \\ \hline
    \ZSFrowColor
    Instabil: & Der Realteil von mindestens einem Eigenwert ist positiv \\
\end{fulltablebox}
